\scanpage{42}\mspage{6}
Pour le voir, on va montrer comment la
carte $\mathcal X$ peut être reconstruite à partir de
l'arrangement 1-dimensionnel orienté. Par la
théorie générale des cartes, il revient au même,
\deleted{\unsure{de}} comme les sommets de $K$ sont d'ordre
$\geq4$ (\deleted{\unsure{$\geq$}}à fortiori $\geq3$), de construire, à partir
de $(K,(D_i)_{i\in I},\omega)$, :

1\textsuperscript{o}) Pour \deleted{tout} $s\in S$ (ens. des sommets de $K$) \deleted{\unsure{le}}
l'ens. \[\inserted{\vec{\unsure{Z|\mathcal Z|\mathfrak Z}}_s\unsure{\in|\subset}\deleted{A}{\unsure{\mathfrak S|\mathscr S|\mathcal C}}_{\vec A_s}}\tag*{(19)}\] des deux ordres circulaires \inserted{(inverses l'un de l'autre)}
canoniques sur l'ens. \inserted{$\vec A_s$} des \deleted{arêtes}\inserted{arcs} de $K$ d'origine
$s$, qu'on peut identifier \deleted{au} \unsure{sic|à} : l'ens
$\vec I_s$ des ps. droites \uline{orientées} \inserted{$\vec D_i$} passant par $s$.

2\textsuperscript{o}) Pour \deleted{\unsure{deux sommets $s,s'$|toute arête orientée}} \deleted{tout} arc $\vec a\in\vec A$ de $K$, d'origine
$s$ et d'extrémité $s'$, la bijection canonique
\[\psi_{\vec a}:\vec Z_s\xrightarrow{\sim}\vec Z_{s'}\tag*{(20)}\]
(qui donne lieu à \unsure{la|l'} \deleted{\unsure{ég}}formule
\[\psi_{-\vec a}=(\psi_{\vec a})^{-1}\tag*{(21)}\]

Quand je dis ``canonique'', j'entends :
provenant de la donnée de la carte $\mathcal X$.
Inversement, si on a pu, à partir de $(K,(D_i)_{i\in I},\omega)$,
construire intrinsèqu\textsuperscript{t} des données
(19) et (20) satisfaisant (21), on trouve une
immersion cellulaire combinatoire de
$K$ dans un $\mathcal X$ (peut être pas le plan projectif !).
% archive sheet number: 41
